\section{System Model and Threat Assumptions}
\label{sec:03}
\label{sec:system-model}

\subsection{System scope and participating entities}
\label{subsec:system-entities}

The earlier voting architecture combined off-chain fingerprint and facial
verification with blockchain-based voter registration and election
operations \citep{rifat2026biometric}. The present study retains this
identity-verification foundation but changes ballot handling after
authentication. In the private-voting pathway, the biometric authorization
request excludes the selected candidate. The voter-facing web client
constructs an encrypted ballot and proves its validity. The voting contract
then accumulates encrypted candidate-slot totals without publishing the
selected candidate. The biometric recognition results belong to the earlier
study. They are not measurements of the cryptographic voting procedure
examined here.

The system involves the voter-facing web client, biometric authentication
service, credential authorization service, election administrator, smart
contracts, public blockchain observers, and five tally trustees. The
administrator configures election parameters, controls lifecycle transitions,
and submits the final result transaction. The biometric authentication service
identifies the voter. The credential authorization service manages registered
voting credentials and issues authorization data. The web client's privacy
worker constructs encrypted ballots and zero-knowledge proofs. The trustees
provide partial-decryption evidence for the result-publication procedure.
These roles have separate security responsibilities. Their separation does
not establish independent organizational control in the prototype.

Registration uses a Semaphore-based group and election anchor. The
private-voting pathway instead uses a custom Poseidon commitment tree and
the \texttt{VoteChainDirect} contract. An accepted private ballot is checked
against the custom eligibility root configured for \texttt{VoteChainDirect}.
The two commitment structures therefore serve different purposes. They do
not form a single eligibility proof.

% Implementation evidence at Git revision 17c7453c50793643b590fb5beac6ad77b9aa53ce:
% src/ZKLocalVoting.sol and contracts/wp4-foundry/src/VoteChainWP4.sol
% votechain-central-server/services/zk/wp5PrivacyService.js
% votechain-frontend-ts/src/components/Voter/VoteCasting.tsx

\subsection{Authorization, membership witness, and ballot submission}
\label{subsec:system-authorization-path}

Before voting, the custom eligibility registry associates each enrolled voter
with a Poseidon identity commitment. This commitment is derived from a private
credential secret. For each election, the election management service
constructs a Merkle tree from registered custom commitments. Its root is
configured in the private-voting contract. The contract uses this public
root to verify ballot eligibility. The credential secret is not published
on-chain.
The voter first selects a candidate in the web client. The selection stays
within the voting interface while biometric authorization takes place. A
separate face or fingerprint request carries an ephemeral credential public
key. The dedicated authorization endpoints reject additional request-body
fields, including the candidate identifier.

After a successful biometric match, the credential authorization service
identifies the enrolled voter and checks the custom credential registration.
It retrieves the Merkle siblings and path directions associated with that
voter's commitment. The service also checks the resulting eligibility root
against the root configured in the active \texttt{VoteChainDirect} contract.
The credential authorization service recovers the stored credential secret.
It then protects the secret for delivery to the requesting web client. This
process uses an ephemeral P-256 key agreement, HKDF-SHA-256 key derivation,
and AES-256-GCM encryption. The authorization response contains the protected
credential envelope, Merkle membership witness, election identifier, and an
expiry time.

The browser privacy worker holds the corresponding ephemeral private key and
opens the credential envelope locally. It recomputes the identity commitment
and Merkle root. The worker then checks the election identifier, candidate
manifest, encryption public key, and election-context binding against the
public election configuration. Any mismatch stops ballot preparation before
proof generation.
Once these checks pass, the worker constructs a one-hot ballot vector for the
locally selected candidate. The vector uses a fixed 16-slot representation.
The worker encrypts each slot under the election's BabyJub public key using
exponential ElGamal. It also derives the election-scoped nullifier.

The worker next generates a Groth16 proof. This proof binds the submitted
ciphertext to the eligibility and ballot-validity constraints. The proof
system follows the witness-hiding construction of Groth
\citep{groth2016snark}. Before submission, the worker checks the public
signals and verifies the generated proof locally.
The transaction interface receives the proof, nullifier, and ciphertext
coordinates. It does not receive the plaintext candidate identifier. These
values are submitted through an Ethereum-compatible vote-submission wallet
to \texttt{VoteChainDirect.castVote}. The contract reconstructs the configured
public inputs and rejects previously used nullifiers. It then verifies the
Groth16 proof. Only an accepted ballot is incorporated into the encrypted
tally.

% Implementation evidence:
% votechain-central-server/middleware/wp5AuthorizationSecurity.js:142-195
% votechain-central-server/controllers/wp5Controller.js:206-240
% votechain-central-server/controllers/fingerprintController.js:131-179
% votechain-central-server/services/zk/wp5PrivacyService.js:482-578
% votechain-frontend-ts/src/privacy/wp5Privacy.worker.ts:680-814,940-1147
% votechain-frontend-ts/src/components/Voter/VoteCasting.tsx:253-315
% contracts/wp4-foundry/src/VoteChainWP4.sol:825-976

\subsection{Public information and identity linkage}
\label{subsec:information-boundaries}

Each registered voter has a private voting-credential secret $s$. The election
identifier is represented by $e$. The circuit derives a nullifier specific
to that election:
\begin{equation}
  n_e = \mathsf{Poseidon}(d_{\mathrm{null}},s,e).
  \label{eq:system-nullifier}
\end{equation}
Here, $d_{\mathrm{null}}$ is the configured nullifier domain separator.
Poseidon is an established hash primitive for zero-knowledge circuits
\citep{grassi2021poseidon}. The membership and ballot-selection witnesses
remain private circuit inputs. The public inputs contain the election
configuration, nullifier, and ciphertext coordinates.

A blockchain observer can inspect an accepted private-vote transaction and
its election-specific nullifier. However, the transaction does not disclose
the credential secret or plaintext candidate choice. Candidate
confidentiality depends on ballot encryption, not on the nullifier
calculation.
The credential authorization service has more information than a public
blockchain observer. It performs identity-aware credential release and
associates stored credentials with identified voters. During authorization,
the service can decrypt a stored credential secret to prepare the protected
browser envelope.

A party with access to this decryption capability can derive $n_e$ from $s$
and $e$. The result can be compared with public voting transactions. This
creates a potential link between an identified voter and an accepted ballot
transaction.
Such a link does not reveal the candidate choice by itself. Recovering that
choice would additionally require compromising the client's selection,
defeating the ballot-encryption confidentiality conditions, or obtaining
sufficient trustee decryption authority. Groth16 hides the private witness
from the proof verifier. It cannot conceal a secret from a service that
already possesses it.

The election administrator and credential authorization service have
different capabilities. Administrative authority to configure or finalize an
election does not itself provide the credential-decryption key. Conversely,
the credential service possesses information unavailable through an ordinary
blockchain transaction. Table~\ref{tab:system-visibility} summarizes these
information boundaries.

\begin{table}[t]
\centering
\caption{Information exposure and verification responsibilities in the private-ballot pathway.}
\label{tab:system-visibility}
\small
\begin{tabular}{@{}p{0.23\linewidth}p{0.34\linewidth}p{0.35\linewidth}@{}}
\toprule
\textbf{Entity} & \textbf{Information available} & \textbf{Security consequence} \\
\midrule
Credential authorization and biometric authentication services
& Voter identity, stored credential, eligibility data, and authorization request
& Can authorize without receiving the candidate choice. Access to the credential-decryption key permits computation of the voter's nullifier. \\
\addlinespace
Voter-facing web client and privacy worker
& Candidate choice, released credential, membership path, encryption randomness, and proof witness during preparation
& Protecting the local selection and credential depends on client integrity. \\
\addlinespace
Public blockchain observer
& Election context, transaction metadata, nullifier, proof, ciphertext slots, and encrypted aggregate
& The accepted private-vote call omits the plaintext candidate and civil identifier. \\
\addlinespace
Election administrator
& Configured election state, encrypted aggregate, finalization hash, and submitted totals
& On-chain administrative authority alone does not establish possession of voter credentials or trustee shares. \\
\addlinespace
Tally trustees and transcript auditor
& Public encrypted tally and verification transcript, with each trustee additionally holding its own secret-share material
& At least three valid trustee contributions are required by the verifier. The confidentiality analysis depends on trustee-share custody. \\
\bottomrule
\end{tabular}
\end{table}

\subsection{Ballot acceptance, tally integrity, and threat conditions}
\label{subsec:system-security-objectives}

Ballot acceptance depends on the cryptographic proof and the contract's state
checks. The circuit constrains membership in the configured custom
eligibility tree. It also binds the nullifier to the election and verifies the
election context. Ballot selection must be one-hot within the active candidate
count. Unused candidate slots must contain zero-valued choices, and the
ciphertexts must satisfy the specified construction.
The contract permits ballot submission only during the active voting phase.
It rejects a reused nullifier or a ciphertext that fails canonical
representation checks. The submitted Groth16 proof must also verify against
the configured public inputs. These checks establish ballot validity under
the implemented circuit relation. They also prevent a second accepted
transaction using the same election nullifier.
Civil-identity uniqueness and credential issuance remain responsibilities of
enrollment and authorization. An accepted on-chain proof does not establish
that a live biometric match occurred at the moment of transaction submission.

During active voting, the contract combines encrypted candidate slots from
accepted ballots. Plaintext totals are recovered through a separate
finalization procedure after the election ends. The external verifier first
checks the distributed key-generation transcript. It then validates trustee
partial-decryption evidence for each candidate slot.
At least three valid trustee shares are required before the verifier recovers
the aggregate candidate totals. The election management service compares
the encrypted aggregates and accepted-ballot count in the transcript with
the corresponding on-chain state. After successful verification, the
administrator submits the recovered totals and finalization-transcript hash.
The contract checks the election phase and the previously pinned
trustee-transcript hash. It also enforces candidate padding, per-slot bounds,
and equality between the reported total votes and the accepted-ballot count.

The distinction between external and on-chain verification is important for
auditability. The contract does not verify the trustees' decryption proofs or
independently recover plaintext totals from the encrypted aggregate. Instead,
it records a finalization hash that identifies the transcript used in the
separate verification procedure.
A successful finalization transaction therefore confirms compliance with the
contract's state and arithmetic checks. It does not, by itself, establish
the validity of every trustee partial decryption. An independent auditor must
obtain the full transcript and verify it against the published encrypted
aggregate and pinned hashes.

% Implementation evidence:
% circuits/lib/VoteChainVoteCore.circom:87-214,305-487
% crypto/src/wp6-threshold.mjs:2307-2595
% votechain-central-server/services/elections/electionFinalizationService.js:994-1257
% contracts/wp4-foundry/src/VoteChainWP4.sol:532-616,692-749,919-976

The threat model includes public observers who can inspect contract state,
transactions, senders, timestamps, and proofs. It also considers ballot
submitters who attempt malformed or repeated votes. The cryptographic
validity relation and contract checks address the specified invalid-ballot
and nullifier-reuse cases.
These mechanisms do not prevent every form of information exposure. A
compromised browser may reveal the selected candidate or credential. A
malicious credential custodian may issue credentials without authorization.
External transaction metadata may also support voter correlation.

Ballot confidentiality depends on the election encryption key and the
trustee-share custody arrangements. Three collaborating trustees with
sufficient valid shares could recover decryption capability. This could allow
them to decrypt individual public ballots as well as the aggregate. The
configured three-out-of-five threshold does not itself establish independent
organizational custody of the shares.
These conditions define the limits of the evaluated system. The available
implementation and experiments do not establish coercion resistance or
receipt-freeness. The local-chain research environment also does not
demonstrate production election deployment. The following section explains
how these security responsibilities are connected through the implemented
voting architecture.